\chapter{SWMF and external background coupling}\label{ch:coupling}
\section{Component lifecycle}
SWMF builds model libraries into one framework executable. AMPS supplies PT
(particle tracker) and PC (particle-in-cell) wrappers in \path{srcInterface}.
The control module owns component registration, processor layout, sessions,
execution/coupling schedules and framework restart requests.
An AMPS standalone executable does not pass through the same outer driver.

The wrapper actions include VERSION, MPI, READ, CHECK, STDOUT, FILEOUT and
GRID, followed by session initialization, run, restart and finalize methods.
READ can occur multiple times before CHECK. Fields and mesh descriptions
must not be published before their initialization stage. The current AMPS
wrappers contain partly stubbed checking/output-routing actions; this manual
does not promote them to a validated full implementation.

\section{Communicators and processor layout}
Use the communicator supplied by the framework. The Fortran handle is
converted with \code{MPI_Comm_f2c}; collective calls must involve exactly
that component's participating ranks. Overlapping component processor sets
execute sequentially on shared ranks. Disjoint sets permit concurrent
evolution, subject to data dependencies. Thread support and actual affinity
must be checked separately from processor counts.

A schematic two-component layout, in a SWMF \path{PARAM.in}, is:
\begin{lstlisting}
#COMPONENTMAP
GM  0  1  1
PT  2  3  1
\end{lstlisting}
It requires four ranks and compatible compiled components. This is a layout
fragment, not a complete runnable physical case.

\section{Parameter ownership and schedules}
SWMF compile-time model selection uses its own \code{Config.pl -v=...}.
Its \code{#BEGIN_COMP PT} block carries component parameters; standalone
AMPS generic runtime input is a different grammar.
Always validate against the actual SWMF \path{PARAM.XML} and component XML:
some historical prose examples reverse positional arguments.
For example, a one-way physical-time cadence uses:
\begin{lstlisting}
#COUPLE1
GM          NameSource
PT          NameTarget
-1          DnCouple
10.0        DtCouple
\end{lstlisting}
This is a 10-s schedule, not a physical justification for that interval.
Time-accurate components must respect supplied synchronization limits.
Tight coupling additionally requires a supported master/slave timestep
contract. Permitting a component to step past a coupling epoch changes
temporal accuracy and must be declared.

Three clocks matter: iterations since process start, accumulated steps
including restart history, and physical simulation time. Returning a limit
to the scheduler without advancing a physical state to that epoch is not
evidence of synchronization. Test this explicitly for nonintegral ratios of
the AMPS step to coupling interval and for first-data readiness.

\section{Exchange contracts}
Each exchanged quantity needs a species/fluid mapping, units, frame,
location (cell centre, corner, face or field-line vertex), timestamp and
validity. Transform vectors, velocities and electric fields with the
correct frame convention; rotating-frame velocity is not just a rotation
of an inertial vector. Keep magnetic interpolation compatible with any
required divergence or flux constraint.

Exchange a coherent state generation rather than mixing fresh density with
old velocity. Verify owner and received-ghost values independently. A
successful send count is not proof of correct physical data. Negative tests
should reject missing variables, nonfinite values, stale epochs, mismatched
layout and unsupported fluid counts before particles consume them.

\section{File providers, SPICE and empirical models}
Offline backgrounds need a versioned file format, dimensional identities,
frame/epoch and explicit missing-domain policy. CDF/Kameleon, BATSRUS/BATL,
ARMS, CCMC and SWMF are separate optional interfaces with different fixtures.
Finding a library on disk does not qualify its adapter.
SPICE supplies geometry, not plasma. Record exact kernel versions, load order,
coverage and target/frame/aberration choices. Kernel readiness requires real
state/orientation queries for the requested epoch.

\section{Coupled testing and isolation}
SWMF's legacy test Makefile serializes its top-level recipes and suppresses
some command failures while retaining \path{.diff} artifacts. A wrapper runner
must inspect expected outputs and classify missing evidence as ERROR.
Never concurrently reconfigure the same SWMF/AMPS build tree for independent
tests. Run isolated runtime cases from frozen executables after serialized
configuration/build, keeping each result separate before aggregation.

The AMPS global SWMF provider case remains SKIP without an independently
registered exchange fixture. Synthetic ABI or replay tests are valuable but
do not replace a live component-communicator and exchanged-field test.
\source{srcInterface/PT_wrapper.f90; srcInterface/PC_wrapper.f90; srcInterface/amps2swmf.cpp}
